\documentclass[11pt]{article}
\usepackage[margin=1in]{geometry}
\usepackage[T1]{fontenc}
\usepackage{lmodern,microtype,amsmath,amssymb,booktabs,array}
\usepackage[colorlinks=true,linkcolor=blue,urlcolor=blue]{hyperref}
\setlength{\parindent}{0pt}
\setlength{\parskip}{5pt}
\newcommand{\reading}[2]{\href{#1}{#2}}
\begin{document}
\begin{center}
{\LARGE Efficiency across generations}\par
\vspace{0.5em}
A reading guide for the housing model \quad September 5, 2026
\end{center}

\section*{1. Whose welfare enters the comparison?}

An overlapping-generations economy contains households with finite lives,
but an endless succession of households. Two periods of life do not mean
that the economy contains only two generations in total.

For a full Pareto comparison, include the initial old, the initial young,
and every future cohort. Compare each household's lifetime utility under the
reform with its utility along the original equilibrium path, starting from
the same inherited situation. For the initial old, their past is fixed, so
only their remaining life can change. With heterogeneous households, an
average improvement within a cohort is insufficient.

For this discussion only, let $\Delta U_i$ denote the change in the lifetime
utility of person $i$; this does not replace the model's value-function notation.
A Pareto improvement requires
\[
\Delta U_i\geq 0\quad\text{for every person in every cohort},
\qquad \Delta U_j>0\quad\text{for at least one person }j.
\]
There is no need to sum these utilities or choose social discount weights.
Those choices arise when specifying a social objective that may trade one
person's loss against another's gain.

Balasko and Shell (1980, p.~286) give a precise definition for an infinite
sequence of generations. They distinguish full Pareto optimality from the
weaker requirement of admitting no improvement that changes only finitely
many households' allocations. Their terminology for this weaker requirement
is specific to that paper. The economic distinction is what matters here.

\textbf{Including a cohort does not require changing its allocation.}
If a feasible reform leaves its entire lifetime allocation unchanged, its
welfare comparison is already settled. A finite construction can therefore
prove inefficiency of the full economy. But showing that no finite
construction works does not exclude an improvement involving infinitely many
cohorts.

\section*{2. A small example}

Consider a separate teaching economy with one young and one old person at
each date. Each person receives two units of goods when young and one when
old. Goods cannot be stored, there are no privately traded assets, and utility
is $\log c^Y+\log c^O$. Without transfers, consumption is $(2,1)$.

Suppose every young person gives half a unit to the current old person, at
every date. Each young person's lifetime consumption becomes $(1.5,1.5)$:
\[
\log 1.5+\log 1.5=\log 2.25>\log 2=\log 2+\log 1.
\]
The initial old person also gains, consuming $1.5$ instead of $1$. Consumption
still totals three units at every date. Everyone gains over their own life,
although every young person initially gives something up.

If the arrangement stops, the last young person who pays does not receive
the later transfer. The absence of a last generation matters. This original
example illustrates Samuelson's intergenerational-transfer intuition; it is
not a specification of our housing economy.

\newpage
\section*{3. Three questions in the literature}

\textbf{Can competition be inefficient simply because generations overlap?}
Samuelson (1958) develops the consumption-loan economy and the role of transfers
between successive generations. Balasko and Shell (1980) make the welfare
distinctions precise: in their frictionless exchange economy, competitive
allocations cannot be improved by changing only finitely many households,
yet need not be fully Pareto optimal. They allow competitive trade in dated
goods. Thus the classical OLG issue cannot be reduced to the usual statement
that an insurance market is missing.

\textbf{Can an economy accumulate too much capital?}
Diamond (1965, sections~3--4) extends OLG analysis to production. Beyond the
Golden Rule capital stock, reducing capital accumulation can release goods
for consumption and permit a Pareto improvement. In the basic deterministic
model, the familiar comparison between the return on capital and population
growth diagnoses this possibility. Diamond also distinguishes comparing
stationary allocations from finding an improvement that respects initial
conditions.

Our illustrative housing model has a fixed housing stock and an externally
given bond price. We have not identified a Diamond overaccumulation mechanism.
Its housing-allocation question must be established from its own preferences,
financial restrictions, and resource constraints. A comparison of interest
and growth rates would not establish our proposed result.

\textbf{How can we separate efficiency from redistribution along a transition?}
Auerbach and Kotlikoff (1987, chapter~5, section~B.3, pp.~62--64) introduce a
hypothetical \emph{Lump Sum Redistribution Authority}. It compensates cohorts
while the reform and compensation are jointly evaluated in general
equilibrium. The authority balances its present-value budget, with assets or
debt allowed along the path. Transfers go to existing cohorts initially and
to later cohorts on entry. This is an efficiency-measurement device, not a
requirement to propose a practical institution.

That is a useful precedent for our exercise. It does not establish that our
particular transfers are feasible: cash before a housing purchase, future
tax enforcement, and the treatment of initial owners remain model-specific
assumptions to assess.

\section*{4. Full-path welfare and steady-state welfare}

These comparisons answer different questions:

\begin{center}
\begin{tabular}{@{}p{0.26\textwidth}p{0.67\textwidth}@{}}
\toprule
Comparison & What must be established \\
\midrule
Pareto improvement & No initial or future household loses; someone gains. \\[5pt]
Steady-state welfare & A specified household is better off at the new stationary allocation. Transitional cohorts still require evaluation. \\[5pt]
Weighted welfare & Gains and losses are aggregated using an explicit social objective and a well-defined treatment of the infinite horizon. \\
\bottomrule
\end{tabular}
\end{center}

For example, a reform might raise the lifetime utility of every cohort born
after year 20 while harming some earlier cohorts. That is a long-run welfare
gain. The uncompensated reform is not a Pareto improvement. Whether feasible
compensation can remove the losses is a further question.

\newpage
\section*{5. What this means for our planner}

\textbf{The welfare comparison and the planner's powers are separate choices.}
We can evaluate every generation while limiting the planner to transfers at
one date. We can also evaluate every generation while allowing a committed
transfer sequence. The latter gives the planner more ways to rearrange
resources; it does not change the meaning of ``no household loses.''

Bisin's \emph{General Equilibrium Theory} lecture notes (2014, section~3.2.4)
are a useful introduction to constrained efficiency. They specify the
planner's instruments and let subsequent markets clear at endogenous prices.
The incomplete-asset-market literature, including Geanakoplos and Polemarchakis
(1986), provides conditions under which asset redistribution can improve
welfare after spot prices adjust. These results depend on their market and
heterogeneity assumptions; they are not an automatic theorem for our
deterministic housing OLG model.

A one-time gift is one possible restriction, not the default definition of
constrained efficiency. Our local obstruction to that instrument does not
settle broader transfer classes. The committed-transfer result remains
conditional on its financing and ownership permissions. The literature helps
frame these choices; each of our proofs still needs its own verification.

\textbf{Fertility introduces a further question.} Holding individual fertility
and cohort sizes fixed lets us compare the same people. When fertility
changes, some people exist in one allocation but not the other. Golosov,
Jones, and Tertilt (2007, section~3.1) study alternative efficiency definitions
for this setting. Their A-efficiency compares people alive in both allocations
over the entire history; it does not mean only people alive at the reform date.
Their P-efficiency includes potential people and requires specifying welfare
when they are not born. A larger population alone is not a welfare ranking.

Our housing-efficiency claim and the fertility prediction can remain separate.
Predicting a larger population does not itself require a welfare ranking
across populations.

\section*{6. A short reading route}

For our immediate question, begin with Auerbach and Kotlikoff, pp.~62--64.
Use Samuelson and Balasko--Shell for foundations, Diamond for capital accumulation,
and Bisin for constrained efficiency. Defer population welfare until we need it.
The links below lead to primary texts or publication records.

{\small
\begin{itemize}
\setlength{\itemsep}{0pt}
\item \reading{https://ludwigbc.com/wp-content/uploads/2013/08/samuelson_overlapping.pdf}{Samuelson (1958)}, ``An Exact Consumption-Loan Model of Interest with or without the Social Contrivance of Money,'' \emph{JPE} 66, 467--482.
\item \reading{https://karlshell.com/wp-content/uploads/2015/03/OGM_1.pdf}{Balasko and Shell (1980)}, ``The Overlapping-Generations Model, I: The Case of Pure Exchange without Money,'' \emph{JET} 23, 281--306. Definitions: p.~286.
\item \reading{https://assets.aeaweb.org/asset-server/files/9443.pdf}{Diamond (1965)}, ``National Debt in a Neoclassical Growth Model,'' \emph{AER} 55, 1126--1150. Begin with pp.~1128--1129.
\item \reading{https://kotlikoff.net/wp-content/uploads/2019/04/Dynamic-Fiscal-Policy.pdf}{Auerbach and Kotlikoff (1987)}, \emph{Dynamic Fiscal Policy}. Chapter~5, section~B.3, pp.~62--64.
\item \reading{https://bpb-us-e1.wpmucdn.com/wp.nyu.edu/dist/c/16384/files/2019/11/Lecture-notes-Oct2014.pdf}{Bisin (2014)}, \emph{General Equilibrium Theory}, lecture notes, section~3.2.4. Technical background: \reading{https://doi.org/10.1017/CBO9780511983566.007}{Geanakoplos and Polemarchakis (1986)}, ``Existence, Regularity, and Constrained Suboptimality of Competitive Allocations When the Asset Market Is Incomplete.''
\item \reading{https://jenni.uchicago.edu/Fertility/Golosov_Jones_Tertilt_2007_Econometrica_v75_n4.pdf}{Golosov, Jones, and Tertilt (2007)}, ``Efficiency with Endogenous Population Growth,'' \emph{Econometrica} 75, 1039--1071. Section~3.1.
\end{itemize}
}
\end{document}
